\section{Incomplete Decoupling and Neff}\label{sec:incomplete_decoupling}
%======================================================================

\subsection{The Physical Problem}

In the standard FLRW picture, neutrino decoupling at $T \sim 1$--$3\;\mathrm{MeV}$ is \emph{incomplete}: during $e^+e^-$ annihilation ($T \lesssim m_e$), residual weak interactions transfer part of the electromagnetic entropy to neutrinos.  This raises the total neutrino energy density above the instantaneous-decoupling value, giving the Standard-Model benchmark $\Neff = 3.044$ instead of $3.000$ \cite{Mangano2005, deSalas2016, Froustey2020, Bennett2020}.  Physically, the correction combines an overall reheating effect with a momentum-dependent spectral distortion, since high-energy neutrinos continue to scatter longer than low-energy ones.

For the present report, however, the key scope statement is different from the textbook one: \emph{the historical exact-characteristic Type~I figures do not rely on a full quantum-kinetic or fully collision-coupled anisotropic Boltzmann treatment of incomplete decoupling}.  Instead, the reference figures rest on a collisionless characteristic-transport core, while incomplete-decoupling effects enter through a tiered approximation hierarchy.  The repository retains historical tier~2 and no-Teff tier~3 JAX surfaces, but these do not amount to a full anisotropic collision-coupled scattering closure.

\paragraph{Current backend authority.} Rust AOT is the active implementation and repeated-run target.  SciPy/BDF remains the temporary number-of-record until the Rust endpoint authority gate passes.  JAX is frozen as a historical component-level parity/automatic-differentiation/Jacobian oracle; it is not an active forward-development, endpoint, runtime-promotion, or production authority.

\paragraph{Current F09B Rust slice.} A crate-private Rust endpoint now evolves
three thermal temperatures with a finite-vacuum-electron-mass, FD/Pauli,
tree-level electron/positron energy-transfer moment.  The same source is
consumed by the corrected weak leg and selected 31-reaction network.  At the
QED-off FLRW anchor $T_\gamma=0.005\,\mathrm{MeV}$, an independently written
adaptive-EOS/CM-integral solve gives
$N=7.93669848685$, $T_{\nu_e}=0.00358318527739\,\mathrm{MeV}$,
$T_{\nu_x}=0.00357679030100\,\mathrm{MeV}$, and the thermal readout
$N_{\rm eff}=3.03409326175$.  DOP853 and Radau agree to
$5.4\times10^{-12}$ in $N$ and below $3.0\times10^{-13}\,\mathrm{MeV}$ in
the neutrino temperatures before the Rust comparison.  This number is a
regression result for the declared thermal tree-level model, not a precision
Standard-Model prediction: spectral transport, neutrino--neutrino reactions,
oscillations, collision QED corrections, and an independently validated
QED-on nuclear endpoint are absent.

\subsection{Tiered Treatment in RABBIT}

The safest description of the current implementation is therefore not ``full treatment'' but \emph{tiered treatment of incomplete decoupling}.  The code contains several levels of thermodynamic and collision-aware sophistication, and the current repository surface is broader than the historical exact-characteristic figures alone.  Table~\ref{tab:decoupling_scope} summarises the intended interpretation.

\begin{table}[t]
\centering
\small
\caption{Historical scope of the incomplete-decoupling hierarchy. ``Used for the quoted report figures?'' refers specifically to the historical exact-characteristic reference results discussed in this report; implementation inventory does not imply current backend authority.}\label{tab:decoupling_scope}
\begin{tabularx}{\textwidth}{p{1.35cm}X>{\centering\arraybackslash}p{1.35cm}>{\centering\arraybackslash}p{1.8cm}p{3.0cm}}
\toprule
Level & Description & Impl.? & Used in quoted figures? & Claim scope \\
\midrule
Tier~1 & Parametric / effective thermodynamic closure; instantaneous-decoupling baseline plus production helper thermodynamics. & Yes & Yes & Historical reference baseline; the retained JAX implementation is a frozen component oracle only. \\
Tier~2 & Coupled three-temperature incomplete-decoupling extension with momentum-averaged $\nu$--$e$ energy transfer and QED-plasma feedback. & Yes & No & Historical JAX component surface; not the F09 collision or endpoint authority. \\
Tier~3 & No-Teff three-temperature weak-budget ladder with CL0--CL3 support. & Yes & No & Historical JAX component surface; not a full anisotropic collision-coupled Boltzmann solver. \\
Tier~4 & Tier~3 plus flavour oscillations and the full SM target. & Deferred & No & Explicitly out of scope. \\
\bottomrule
\end{tabularx}
\normalsize
\end{table}

This hierarchy should be read with three distinct labels in mind.
\begin{enumerate}
\item \textbf{Reference-result claim}: the historical exact-characteristic figures and CL2 profile constraint quoted elsewhere in this report.
\item \textbf{Frozen component-oracle claim}: the retained JAX tier~1--3 surfaces may support local parity, AD, and Jacobian checks only; they do not establish forward-solver or runtime authority.  SciPy/BDF is the temporary number-of-record, not a Rust promotion result.
\item \textbf{Out-of-scope claim}: full QKE, full anisotropic collision integrals, species-resolved end-to-end scatter on exact rays, and complete $\nu$--$\nu$ plus $\nu$--$e$ angular closure.
\end{enumerate}

With that distinction in place, the phrase ``incomplete decoupling'' in the remainder of this document should always be read as a \emph{controlled approximation hierarchy} rather than as a claim that the full anisotropic scattering problem has already been solved end to end.

\subsection{Entropy Conservation with Collision Source Term}

The fundamental thermodynamic equation governing photon-temperature evolution is entropy conservation of the $\gamma+e^\pm$ plasma, augmented by a source term describing neutrino--electron energy exchange.  Let $\mathcal S = s_{\rm plasma}a^3$ be the comoving electromagnetic entropy, with $s_{\rm plasma}=(\rho_{\rm plasma}+p_{\rm plasma})/T_\gamma$.  Then
\begin{equation}
\frac{\dd}{\dd t}(s_{\rm plasma}a^3) = -\frac{1}{T_\gamma}\,\frac{\dd Q_{\nu\to{\rm plasma}}}{\dd t}\,a^3\,,\label{eq:entropy_full}
\end{equation}
where $\dd Q_{\nu\to{\rm plasma}}/\dd t$ is positive when neutrinos lose energy to the plasma.

Writing $s_{\rm plasma}=S(T_\gamma)T_\gamma^3$ and differentiating with respect to $N=\ln a$ gives
\begin{equation}
\frac{\dd}{\dd N}\bigl[S(T_\gamma)T_\gamma^3\bigr] = \bigl(3S+T_\gamma S'\bigr)T_\gamma^2\frac{\dd T_\gamma}{\dd N}+3ST_\gamma^3\,.
\end{equation}
Solving for the temperature derivative yields the thermodynamic master equation
\begin{equation}
\bx{\frac{\dd T_\gamma}{\dd N} = -\frac{T_\gamma}{1 + T_\gamma S'/(3S)}\left(1 - \frac{\dd Q}{T_\gamma^4\,(\dd\rho/\dd T)_\gamma}\right)\,.}\label{eq:dTgamma_dN}
\end{equation}
In the historical hierarchy, Tier~1 supplied the source term through an effective closure, while Tier~2 used a partially live collision-aware module.  This describes the approximation structure only; it does not promote either historical JAX surface to current collision or endpoint authority.

At Tier~2+, the energy-transfer functional is obtained from the collision operator:
\begin{equation}
\frac{\dd Q_{\nu\to\rm plasma}}{\dd t} = -\sum_\alpha \int \frac{\dd^3\mathbf p}{(2\pi)^3}\,E\,C_\alpha[f_\alpha] \, ,
\end{equation}
and, using the momentum variable $y=p/T_\nu$,
\begin{equation}
\frac{\dd Q}{\dd N} = -\frac{T_\nu^4}{2\pi^2 H}\sum_\alpha\int_0^\infty \dd y\; y^3 C_\alpha[f_\alpha](y)\,.
\end{equation}
These expressions are kept because they are the natural bridge between the thermodynamic equation and the Tier~2 extension, even though the fully anisotropic evaluation of $C_\alpha$ remains outside current endpoint authority.

\subsection{Momentum-Dependent Decoupling}

\begin{figure}[t]
\centering
\includegraphics[width=\textwidth]{fig_ped02_collision_rate_profile.png}
\caption{Neutrino collision rates.  (a)~$\Gamma/H$ vs photon temperature for $\nu_e$ (blue, CC+NC coupling) and $\nu_x$ (red, NC only).  The decoupling condition $\Gamma/H = 1$ is reached at $T \approx 1.5\;\mathrm{MeV}$ for $\nu_e$ and $T \approx 3.5\;\mathrm{MeV}$ for $\nu_x$.  (b)~Momentum-dependent decoupling temperature $T_{\rm dec}(y) \propto y^{-1/3}$: high-momentum neutrinos decouple later and receive more heating from $e^+e^-$ annihilation.}\label{fig:collision_rate}
\end{figure}

The collision rate for a neutrino with dimensionless momentum label $y$ scales approximately as $\Gamma(y) \propto G_F^2T^5 y$, so high-momentum neutrinos remain coupled longer than low-momentum neutrinos.  A useful estimate of the decoupling temperature is
\begin{equation}
T_{\rm dec}(y) \approx T_{{\rm dec},0}\left(\frac{y}{3.15}\right)^{-1/3},
\end{equation}
with $T_{{\rm dec},0}\approx 1.5\;\mathrm{MeV}$ for $\nu_e$ near the thermal peak.  The consequence is a momentum-dependent reheating pattern: low-$y$ modes decouple early and miss most of the $e^+e^-$ entropy release, while high-$y$ modes remain coupled longer and are heated more strongly.

To display this effect without confusing the dimensionless transport label $y=p/T_\nu$ with a physical temperature, it is safer to define a \emph{diagnostic effective temperature at fixed physical momentum} $p$ by
\begin{equation}
T_{\rm eff}(p) \equiv \frac{p}{\ln\!\bigl(1/f(p)-1\bigr)}\,.
\end{equation}
For a perfect Fermi--Dirac spectrum, $T_{\rm eff}(p)=T_\nu$ is independent of $p$; after incomplete decoupling, $T_{\rm eff}(p)$ rises at high $p$, signalling spectral hardening.  In the present report this quantity is used only as a diagnostic descriptor of the Tier~2 extension, not as a separate production dynamical variable.

\begin{figure}[t]
\centering
\includegraphics[width=\textwidth]{fig_ped03_spectral_distortion.png}
\caption{Neutrino spectral distortion from incomplete decoupling (schematic, based on Mangano et~al.\ 2005).  (a)~Distribution functions: the distorted $f_{\nu_e}$ (dashed) shows excess at high $y$ relative to equilibrium (solid).  (b)~Fractional distortion $\delta f/f_{\rm eq}$: positive at $y > 3$ and negative at $y < 3$.  (c)~A diagnostic effective-temperature representation showing the corresponding hardening trend.}\label{fig:spectral_dist}
\end{figure}

\subsection{Connection to Anisotropic BBN}

In anisotropic Bianchi cosmologies, incomplete decoupling acquires two distinct layers of relevance.  First, even at the level of the production baseline, the decoupling condition $\Gamma/H$ is altered by the shear-enhanced Hubble rate
\begin{equation}
H=H_{\rm FLRW}/\sqrt{1-\Sig^2}
\end{equation}
so decoupling occurs earlier when $\Sig>0$.  This is the leading anisotropic effect relevant for the present BBN forward model.

Second, once one goes beyond the production baseline, the neutrino distribution ceases to be both isotropic and exactly thermal.  A fully faithful treatment would then require species-, momentum-, and angle-resolved scattering on top of exact ray transport.  That problem is \emph{not} solved here end to end.  The historical Tier~2 modules remain available for bounded component checks, without establishing a complete anisotropic scattering closure or a validated BBN correction.

The report retains Tier~2 for provenance and bounded component comparison.  It is not a production headline, an active runtime path, or validation of the deferred anisotropic collision-coupled problem.

\subsection{\texorpdfstring{$N_{\rm eff}$}{Neff} in the tiered hierarchy}

The effective number of neutrino species is defined operationally by
\begin{equation}
\Neff \equiv \frac{8}{7}\left(\frac{11}{4}\right)^{4/3}\frac{\rho_\nu}{\rho_\gamma}\,.
\end{equation}
What changes across tiers is not this definition but the approximation used for $\rho_\nu$.

At the instantaneous-decoupling level one has
\begin{equation}
\rho_\nu^{\rm inst} = N_\nu\,\frac{7\pi^2}{120}T_\nu^4,
\qquad T_\nu/T_\gamma=(4/11)^{1/3},
\end{equation}
which gives $\Neff=3.000$ before any reheating correction is added.  A purely parametric helper can then impose the Standard-Model benchmark $\Neff=3.044$ by hand.

In the intended historical Tier~2 formulation, the neutrino energy density would instead be computed from the evolved non-equilibrium spectrum,
\begin{equation}
\rho_\nu = \sum_\alpha \frac{T_\nu^4}{2\pi^2}\int_0^\infty \dd y\; y^3 f_\alpha(y)
\end{equation}
but this expression alone does not validate a numerical $\Neff$ prediction.  Missing $\nu$--$\nu$ scattering, flavour oscillations, and full collision closure prevent a Standard-Model precision claim.  The frozen no-Teff tier~3 JAX surface likewise remains a component oracle rather than a full species- and angle-resolved collision solver.

\subsection{The Full Boltzmann Equation for Incomplete Decoupling}

At Tier~2, the evolution equation for each neutrino species $\alpha$ in the thermal momentum variable $y \equiv p/T_\nu$ is
\begin{equation}
\bx{\frac{\partial f_\alpha(y)}{\partial N} = -\frac{\dd\ln y}{\dd N}\,y\,\frac{\partial f_\alpha}{\partial y} + \frac{C_\alpha[f_\alpha, f_e](y)}{H}\,,}\label{eq:boltzmann_tier2}
\end{equation}
where $C_\alpha$ includes both scattering and pair processes.  In the FLRW limit, there is no angular dependence ($\dd\mu/\dd N = 0$), and the problem reduces to solving $N_q$ coupled ODEs per neutrino species for the distribution values $f_\alpha(y_i)$ at each momentum node.

\paragraph{Current F10A collisionless substrate.} The first Rust spectral
slice fixes the comoving coordinate $q=ap$, with $a=1$ at the initial state,
instead of using a moving physical-momentum grid.  For massless collisionless
neutrinos the Liouville equation is therefore represented exactly as
\begin{equation}
\left.\frac{\dd f(q_i)}{\dd N}\right|_{C=0}=0,
\qquad
n_{\nu\bar\nu}=\frac{e^{-3N}}{\pi^2}\sum_iw_iq_i^2f_i,
\qquad
\rho_{\nu\bar\nu}=\frac{e^{-4N}}{\pi^2}\sum_iw_iq_i^3f_i .
\end{equation}
The distribution energy density enters the Friedmann rate and hence an
evolved cosmic-time state in the photon-temperature endpoint; it is not a
detached moment check.  With 24 Gauss--Laguerre nodes, Rust BDF/Rodas5P reach
$N=7.93804264043/7.93804264041$ and elapsed time
$52796.0351546/52796.0351521\,\mathrm{s}$ at
$T_\gamma=0.005\,\mathrm{MeV}$.  A separately written adaptive-EOS continuum-
FD SciPy solve gives $N=7.93804264044$ and
$52796.0351521\,\mathrm{s}$; DOP853 and Radau agree before the Rust comparison.
Orders 12, 24, and 36 converge to the analytic FD number and energy moments,
and arbitrary occupations preserve $n_\nu a^3$ and $\rho_\nu a^4$.

At the F10A boundary this closed only collisionless isotropic classical
redshifting and the energy-grid state convention; it supplied no collision or
distortion claim.

\paragraph{Current F10B2 electron action.} F10B1 first evaluated the
finite-mass HM electron/positron event action on the same grid for separate
electron and degenerate-heavy flavour-pair occupations.  Conservative
antisymmetric elastic fluxes preserve the quadrature number moment, symmetric
pair fluxes preserve neutrino-minus-antineutrino number, and the analytic
Pauli response supplies the occupation Jacobian.  The collision energy moment
is debited equal-and-opposite from the electromagnetic bath.  The exact common
FD reference is an analytic action null; five-point action/FLRW derivatives and
the combined first law pass.  F10B2 stores the evolved occupations as exact
logits, after a direct-occupation 12-node Rodas run failed raw in the
floating-point tail.  Stable logistic reconstruction and the chain-rule
Jacobian preserve strict open occupations without clipping; unrepresentable
underflow still fails raw.

A separated momentum/event ladder rejects event order four/three and retains
16 momentum nodes with event order six/four.  Its electron/heavy thermal energy
moments are $0.99078/0.99087$ of the independent thermal anchors.  BDF and
Rodas5P reach
$(N,t,\Neff)=(7.9367312298,52681.0360\,\mathrm{s},3.0330775769)$ and
$(7.9367453070,52681.5412\,\mathrm{s},3.0329326538)$ at
$T_\gamma=0.005\,\mathrm{MeV}$.  A same-code grid-20 BDF refinement gives
$\Neff=3.0331966346$, while an explicit-$N$ step ladder bounds direct-support
sensitivity.  Grid-20 Rodas and an independent full spectral solve were not
run.  These remain test-only regression readouts: no precision distortion or
$\Neff$ claim, neutrino--neutrino collision, lepton asymmetry, flavour
coherence, anisotropy, QKE, or public runtime claim follows.

\paragraph{Historical F10C0 GL16 and current F10C1 N48 endpoint.} The FLRW
right-hand side consumes the restricted same-flavour
$\nu_\alpha\nu_\alpha\leftrightarrow\nu_\alpha\nu_\alpha$ action and the
identical antineutrino channel.  Because this elastic self-action conserves
its discrete neutrino energy moment event by event, it introduces no
electromagnetic energy-transfer source.  At the historical F10C0 boundary,
with the 16-node Gauss--Laguerre momentum grid, electron event order six/four,
and self-scattering angular order 12, the private cold-event regressions were
\begin{align*}
(N,t,\Neff)_{\rm BDF}
  &= (7.9367292791,52680.9384\,\mathrm{s},3.0331059589),\\
(N,t,\Neff)_{\rm Rodas5P}
  &= (7.9367451102,52681.5251\,\mathrm{s},3.0329371816).
\end{align*}
The selected linear action was approximately $7.54\%$ below the high-
resolution tagged oracle, so these values remain historical
\textbf{IMPLEMENTED} regressions rather than radial or precision authority.

F10C1 retains the same coefficient, Pauli algebra, and conservative linear
deposition but selects the positive exponential N48 rule
$y=-3\ln(1-t)$.  The initial no-cache N48 run gives
\begin{align*}
(N,t,\Neff)_{\rm BDF}
  &= (7.9367214190,52680.2048\,\mathrm{s},3.0333103242),\\
(N,t,\Neff)_{\rm Rodas5P}
  &= (7.9367373740,52680.8451\,\mathrm{s},3.0331291370).
\end{align*}
The direct N32--N64 ladder decreases from $1.873525\%$ to $0.489704\%$,
and two independent N128 auxiliary integrations bound five declared smooth
nonthermal profiles by $1.9771546\%<2.1\%$.  This is
\textbf{VALIDATED} only as a profile-bounded radial representation; it is not
arbitrary-distribution, continuum, complete-catalogue, precision-$\Neff$, or
independent full-spectral authority.

The initial matched no-cache endpoint took $2181.07\,\mathrm{s}$ with maximum
RSS $51020\,\mathrm{KB}$.  A prepared self-event-geometry candidate took
$1980.52\,\mathrm{s}$, only $9.195\%$ faster and below the programme's
approximately $10\%$ partial-blocker threshold, while RSS rose by $3.277\%$.
Its 112 Rust lines and cache-only test were removed; this negative result is
\textbf{DEPRECATED} evidence, not endpoint promotion.
The retained no-cache tree repeats both solver outputs bitwise in
$2191.66\,\mathrm{s}$ with maximum RSS $50440\,\mathrm{KB}$, respectively
$+0.486\%$ and $-1.137\%$ relative to the initial no-cache run.  The cold
wall is therefore unchanged.

Catalogue row 1 is complete and the same-flavour-identical subset of row 3
executes.  The distinct-heavy remainder of row 3 and every channel in rows 2
and 4--9 remain \textbf{SPECIFIED} but absent, so row 3 is partial and seven
rows are wholly absent.  Anisotropic collision closure, oscillations, QKE,
and public runtime authority do not follow; the measured N48 cold wall remains
the active task before catalogue expansion.  By owner decision the present
programme stops after this collision-coupled isotropic FLRW milestone is
completed and validated; Bianchi/LRS/non-LRS work requires a separate future
instruction.

When anisotropy is present, the full Boltzmann equation along a characteristic ray $j$ becomes
\begin{equation}
\bx{\frac{\partial f_j(y)}{\partial N} = -\left(\frac{\dd\ln y}{\dd N}\right)_{\!j}\,y\,\frac{\partial f_j}{\partial y} + \frac{C[f](y, \mu_j)}{H}\,,}\label{eq:boltzmann_char_coll}
\end{equation}
where $(\dd\ln y/\dd N)_j$ is the gravitational redshift along ray $j$ (determined by the anisotropic geometry), and $C[f](y, \mu_j)$ is the collision integral evaluated at the direction $\mu_j$ of ray $j$.

This equation poses a fundamental challenge: the collision integral $C[f](y, \mu_j)$ requires knowledge of $f$ at \emph{all} directions $\mu$ (through the angular integration in the matrix element), but each characteristic ray only carries information about its own direction $\mu_j$.  The current production report therefore treats the fully anisotropic collision-coupled problem as deferred full phase-space scope rather than as an active reduced bridge.

\subsection{The State Vector for Full Boltzmann}

In the full phase-space ray approach (Paper~II scope), the state per ray becomes $f_j(y_k)$ at $N_q$ momentum nodes, rather than the reduced characteristic variables.  The DOF comparison is:

\begin{center}
\begin{tabular}{lcc}
\toprule
Approach & State representation & Indicative size \\
\midrule
Characteristic reference (explicit) & $(I_j, J_j)$ + shared $S$ & $2N_\mu + 1 = 25$ transport DOF \\
Characteristic JAX compact & shared $S$ with analytic $I_j(S), J_j(S)$ & 13 total DOF (tier~1), 15 total DOF (tier~2) at $N_\mu=12$ \\
Full phase-space ray & $f_j(y_k)$ ($N_q$ values per ray) & $N_\mu \times N_q = 240$ transport DOF \\
\bottomrule
\end{tabular}
\end{center}

The old 37-DOF count refers only to the explicit characteristic reference layout.  The frozen compact JAX characteristic oracle is smaller still, while the hypothetical full phase-space ray path would replace the reduced ray variables by $12 \times 20 = 240$ distribution values $\{f_j(y_k)\}$ and increase the total system size to the few-hundred-DOF range.  This count is an indicative design estimate, not evidence that the full solver is implemented or tractable at endpoint.

\begin{figure}[t]
\centering
\includegraphics[width=\textwidth]{fig_ped07_phase_space_comparison.png}
\caption{Degrees of freedom comparison across transport methods.  (a)~Transport DOF: the linearised PSTF hierarchy requires 960 DOF, while characteristic rays need only 25.  The proposed full phase-space layout uses 240.  (b)~Indicative total system DOF including geometry and network.  (c)~Dense Jacobian entry count scales quadratically; the $\sim70{,}000$-entry estimate does not by itself establish Rosenbrock endpoint tractability.}\label{fig:dof_comparison}
\end{figure}

%======================================================================
